\documentclass[11pt]{article}

\usepackage[a4paper,margin=1in]{geometry}
\usepackage{amsmath,amssymb}
\usepackage[hidelinks]{hyperref}
\usepackage{enumitem}

\emergencystretch=4em
\sloppy

\title{Third Codex Review Findings for\\
\texttt{rational\_cut\_and\_project\_gap\_periods.tex}}
\author{Codex Review Notes}
\date{May 13, 2026}

\begin{document}

\maketitle

\section*{Executive Summary}

I reviewed the current manuscript
\[
\texttt{LaTeX/rational\_cut\_and\_project\_gap\_periods.tex}
\]
together with the current Lean code:
\begin{itemize}[leftmargin=*]
  \item \texttt{Lean/CutAndProject/CutAndProject/Basic.lean}
  \item \texttt{Lean/CutAndProject/CutAndProject/Irrational.lean}
\end{itemize}

The rational main theorem and the set-valued companion theorem look
mathematically solid.  I did not find a blocking issue in the residue
reduction, the multiset period formula, the set-valued period formula, or the
Lean formalisation of the residue-combinatorial core.  The remaining findings
are precision and presentation issues, concentrated mainly in the irrational
section and in the description of the Lean scope.

\section*{Findings}

\begin{enumerate}[label=\textbf{F\arabic*.},leftmargin=*]

\item \textbf{Irrational unboundedness paragraph: make the \(\varepsilon<1\)
condition explicit.}

Around
\[
\texttt{LaTeX/rational\_cut\_and\_project\_gap\_periods.tex:1255--1284},
\]
the argument first uses \(\varepsilon<1\) to force \(p_0\ne 0\) and obtain
the lower bound
\[
|\tilde p(p_0,q_0)| \ge c_0 := 1+a^2-|a|.
\]
Later the chosen smallness condition is written only as
\[
\varepsilon \le
 \frac{c_0\min(a\omega,\omega)}{|M|+2c_0+1}.
\]

\textbf{Recommendation.}  Choose explicitly
\[
\varepsilon \le
\min\!\left(1,\,
 \frac{c_0\min(a\omega,\omega)}{|M|+2c_0+1}\right).
\]
This matches the Lean proof, where the auxiliary parameter is defined using
\(\min 1(\cdots)\) in \texttt{exists\_scaled\_witness}
(\texttt{Irrational.lean:334}).

\item \textbf{Lean scope: keep the wording slightly more conservative.}

The Lean code itself correctly states that the geometric construction is not
formalised from first principles:
\[
\texttt{Basic.lean:505--510}.
\]
The paper is mostly aligned with this, especially in the ``Scope of the
formalization'' subsection.  However, phrases such as ``actual geometric
enumeration'' around
\[
\texttt{LaTeX/rational\_cut\_and\_project\_gap\_periods.tex:1639}
\]
could still be read too strongly.

\textbf{Recommendation.}  Use consistently:
\begin{quote}
``the residue-combinatorial core with the geometric multiplier threaded
through to the final theorems''
\end{quote}
rather than wording suggesting that the full real-geometric
strip-to-residue construction itself is machine-checked.

\item \textbf{Irrational period lift: avoid speaking of a period
\(-\sigma\).}

In the proof of the irrational aperiodicity proposition, around
\[
\texttt{LaTeX/rational\_cut\_and\_project\_gap\_periods.tex:1350--1355},
\]
the text says:
\begin{quote}
``Applying the same argument to the period \(-\sigma\) yields the reverse
inclusion.''
\end{quote}
Strictly speaking, \(-\sigma\) is not a period length of the gap sequence.
It is a translation length in the projected set.

The Lean proof is cleaner: it only produces the forward inclusion
\[
A+v\subseteq A
\]
in
\[
\texttt{period\_lifts\_to\_lattice\_translation}
\quad(\texttt{Irrational.lean:704--709}),
\]
and
\[
\texttt{lattice\_translation\_must\_be\_zero}
\quad(\texttt{Irrational.lean:798--801})
\]
proves that this forward inclusion alone
forces \(v=0\).

\textbf{Recommendation.}  Align the paper proof with Lean: claim only
\(A+v\subseteq A\), and in Step~3 derive the contradiction from this forward
inclusion by choosing an internal point near the appropriate endpoint of the
window.

\item \textbf{Define \(A\) before using it in the irrational preamble.}

The notation \(A\) is used in the explanatory paragraph before
the irrational aperiodicity proposition, for example in \(\tilde p(A)\) and
\((Kp_0,Kq_0)\in A\), before it is formally defined inside the proposition's
proof.

\textbf{Recommendation.}  Either define \(A\) and \(W\) before that paragraph,
or avoid the notation \(A\) until the proof begins.

\item \textbf{Small wording issue in the period-divisibility lemma.}

In the ``Minimal period divides every period'' lemma, around
\[
\texttt{LaTeX/rational\_cut\_and\_project\_gap\_periods.tex:687--690},
\]
the proof calls \(M-q\lambda\) a ``nonnegative integer combination'' of
periods.  It is really a difference of periods.  The conclusion is correct
for a bi-infinite sequence, but the wording is a little loose.

\textbf{Recommendation.}  Say instead that because \(M\) and \(\lambda\) are
periods, so is \(M-q\lambda\) (using negative shifts if needed), and then
apply minimality.

\end{enumerate}

\section*{Checks Performed}

\begin{enumerate}[leftmargin=*]

\item \textbf{Lean build.}

The command
\[
\texttt{lake build}
\]
inside \(\texttt{Lean/CutAndProject}\) completed successfully:
\[
\texttt{Build completed successfully (8251 jobs).}
\]
No occurrences of \texttt{sorry}, \texttt{admit}, or \texttt{axiom} were
found in the checked Lean source files.

\item \textbf{CSV verification.}

Running
\[
\texttt{python3 src/python/verify\_paper.py}
\]
on the three cited CSV files reported:
\[
\begin{array}{rl}
\text{total rows:} & 51{,}729,\\
\text{multiset mismatches:} & 0,\\
\text{set mismatches:} & 0,\\
\text{negative set sentinels:} & 0.
\end{array}
\]

\item \textbf{Paper build.}

Two runs of \texttt{pdflatex} completed successfully.  After the second run
there were no undefined references, no citation warnings, and no overfull
boxes in the searched log output.  The remaining warnings were only underfull
box warnings in the Lean/table area.

\end{enumerate}

\section*{Overall Assessment}

The paper is close to arXiv-ready.  The rational core is coherent:
\[
r=\alpha x-\beta y,\qquad
\tilde p=\beta x+\alpha y,\qquad
c_r\equiv -\alpha\beta^{-1}r \pmod D
\]
lead cleanly to the residue distribution argument, and the minimality proof
via the stabilizer of a proper cyclic interval is convincing.  The set-valued
period theorem also follows cleanly from the same residue distribution.

Before submission, I would make the small precision edits above, especially in
the irrational section and in the Lean-scope wording.  I do not see a
mathematical blocker in the rational multiset or set-valued results.

\end{document}
